\section{The method: what makes the system}
\label{sec:idea}

\subsection{Knowledge Graphs Overview}
\label{sec:idea-pivot}



\label{sec:idea-three}
Solving Financial RAG requires 2 types of prior knowledge. (1) What different financial Concepts mean and how they are interconnected. (2) How are Financial document structured and what information do they contain. To Incorporate this we create a \emph{concept graph} which says what quantities exist and how they relate to each other, and a \emph{source graph} (Section~\ref{sec:idea-source}), which says what a document of a given kind contains. A \emph{bridge} (Section~\ref{sec:idea-bridge}),  brings them together by connecting the financial concepts to the document kinds. The edges represent the prior likelihood of each concept being found in a particular section of a document.

\begin{figure}[t]
\centering
\fitwidth{\begin{tikzpicture}[
  >=Stealth, font=\small,
  base/.style={draw,thick,rounded corners=3pt,align=left,inner sep=6pt,text width=46mm,font=\scriptsize},
  kb/.style={base,fill=cLight},
  title/.style={font=\small\bfseries},
  arr/.style={->,thick,gray!80}
]
\node[kb,draw=cEntity] (cg) at (0,0) {%
 \textcolor{cEntity}{\textbf{Concept graph}} --- \emph{what} and \emph{how related}\\[2pt]
 9{,}121 entities (4{,}474 textbook, 4{,}647 SEC)\\
 5{,}457 relations: 2{,}753 equations, 1{,}384 subtotals, 1{,}320 rules\\
 4{,}983 links entity $\leftrightarrow$ SEC concept, one entity per concept\\[2pt]
 \emph{verified by}: arithmetic over filings; review of definitions\\
 \emph{example}: current ratio $=$ current assets / current liabilities (identity)};
\node[kb,draw=cDoc,right=14mm of cg] (sg) {%
 \textcolor{cDoc}{\textbf{Source graph}} --- \emph{kinds of documents}\\[2pt]
 15 kinds, 150 official parts, 29 period rules, 20 relations, 4 filer classes\\
 closed vocabulary of $\sim$13 period words\\[2pt]
 \emph{verified by}: a quote from the official SEC text per entry; counts over 68{,}177 filings\\
 \emph{example}: 10-K $\triangleright$ Part II $\triangleright$ Item 8 holds the statements for the fiscal year and the prior year end; a 10-Q never holds Q4};
\node[kb,draw=cConnector,below=12mm of $(cg.south)!0.5!(sg.south)$] (br) {%
 \textcolor{cConnector}{\textbf{Bridge}} --- \emph{where each quantity is printed, how often}\\[2pt]
 49{,}677 rows (concept, kind, part, statement or note, periods, count, of)\\
 measured over 32{,}288 filings of 7 kinds from the SEC Financial Statement (and Notes) data sets; observed rows added from the system's own finds; authored rows as priors\\[2pt]
 \emph{example}: effective tax rate $\to$ 10-K, Item 8, income-tax note, 56\% of filings; capex $\to$ cash flow statement, 61\%};
\draw[arr] (cg.south) -- ++(0,-5mm) -| ([xshift=-12mm]br.north);
\draw[arr] (sg.south) -- ++(0,-5mm) -| ([xshift=12mm]br.north);
\node[font=\scriptsize,gray,align=center,below=3mm of br] {The bridge is keyed on the SEC concept (stable) and resolved to the entity through the concept links when opened.};
\end{tikzpicture}}
\caption{The three knowledge bases. The concept graph says what quantities exist and how they relate; the source graph says what a document of a given kind contains and for which periods; the bridge joins them by saying in which part and statement of which kind each quantity is printed, with counts. Each is verified by a different kind of evidence.}
\label{fig:three}
\end{figure}


\subsection{The Concept graph}
\label{sec:idea-graph}

The graph has two types of nodes (Figure~\ref{fig:capm}). (1) \termdef{Entity nodes} which represent financial quantities (the risk-free rate, total current assets, the effective tax rate), each with a name, aliases, type and units, a plain-language definition and a vector representation. (2) \termdef{Connector nodes} are \(n\)-ary relations which represent how different concepts interact with each other. Connector nodes can store:
\begin{enumerate}
    \item an \emph{equation} in canonical \(f(x_1,\dots,x_n)=0\) form where the inputs are the entities involved in the equation. t also stores the assumptions under which the equation holds
    \item \emph{rules} which store textual information on how different entities affect each other
    \item \emph{aggregation} (a subtotal with a signed, weighted edge per member ever observed).
\end{enumerate}
In the graph edges only connect to each other via connectors. A connector may reference both entities and other connectors. Connectors also have a trust level associated with them based on how likely the relation is to hold in real life.

\begin{figure}[t]
\centering
\fitwidth{\begin{tikzpicture}[
  >=Stealth,
  ent/.style={circle,draw=cEntity,fill=cEntity!12,thick,minimum size=9mm,font=\scriptsize,align=center,inner sep=1pt},
  con/.style={rectangle,draw=cConnector,fill=cConnector!15,thick,rounded corners=1pt,minimum width=36mm,minimum height=11mm,font=\scriptsize,align=center,inner sep=3pt},
  lab/.style={font=\scriptsize\ttfamily,fill=white,inner sep=1pt},
  every edge/.style={draw,thick,gray!70}
]
% ---------- CAPM ----------
\node[con] (capm) at (0,0) {\textbf{equation} \\ $E_{R_i} - R_f - \beta_i\,(E_{R_m} - R_f) = 0$ \\ solvable for every symbol};
\node[ent,above left=11mm and 14mm of capm] (eri) {expected\\return\\of asset $i$};
\node[ent,above right=11mm and 14mm of capm] (rf) {risk-free\\rate};
\node[ent,below left=11mm and 14mm of capm] (beta) {beta of\\asset $i$};
\node[ent,below right=11mm and 14mm of capm] (erm) {expected\\market\\return};
\path (eri) edge node[lab,pos=.45]{E\_R\_i} (capm);
\path (rf) edge node[lab,pos=.45]{R\_f} (capm);
\path (beta) edge node[lab,pos=.45]{beta\_i} (capm);
\path (erm) edge node[lab,pos=.45]{E\_R\_m} (capm);
\node[font=\scriptsize,gray,align=center,below=18mm of capm] {$R_f$ appears twice in the body but is \emph{one} binding, one edge.\\ The realised return of asset $i$ is a different entity, which this equation does not touch.};

% ---------- connector referencing connector ----------
\begin{scope}[xshift=88mm]
\node[con] (d1) at (0,-14mm) {\textbf{equation} \texttt{d1\_def} \\ $d_1 - \dfrac{\ln(S/K)+(r+\sigma^2/2)T}{\sigma\sqrt T} = 0$};
\node[con] (bs) at (0,14mm) {\textbf{equation} \texttt{bs\_price} \\ $C - S\,\Phi(d_1) + K e^{-rT}\Phi(d_1-\sigma\sqrt T) = 0$};
\node[ent,left=9mm of d1] (S) {spot\\price $S$};
\node[ent,right=9mm of d1] (sig) {volatility\\$\sigma$};
\node[ent,left=9mm of bs] (C) {call\\price $C$};
\node[ent,right=9mm of bs] (K) {strike\\$K$};
\path (S) edge (d1); \path (sig) edge (d1); \path (C) edge (bs); \path (K) edge (bs);
\path (bs) edge[dashed,cConnector,very thick] node[lab,right=1pt,pos=.5]{d1 $\to$ \texttt{NodeRef(d1\_def)}} (d1);
\node[font=\scriptsize,gray,align=center,below=4mm of d1] {A connector may bind a symbol to another connector (dashed):\\ the graph is bipartite in spirit, not strictly.};
\end{scope}
\end{tikzpicture}}
\caption{The entity--connector graph. Left: CAPM as one equation connector (orange) bound to four entity nodes (blue) by undirected edges carrying the local symbol. The expected returns are entities in their own right; no edge or expression applies $E[\cdot]$. Right: a connector referencing another connector, the $d_1$ definition used inside the Black--Scholes call price. $r$ and $T$ are omitted from the drawing for space.}
\label{fig:capm}
\end{figure}



To compute a target, the \termdef{solver} searches the graph as an AND/OR tree.
A target at a coordinate (company, period, segment) is an \emph{OR-node}: it can
be produced by any one of the connectors that produce it. Each connector is an
\emph{AND-node}: it can be used only when all its other inputs are available, and
each of those inputs becomes a target in turn. The solver explores this tree by
iterative deepening and memoises failed targets so that they are not tried twice. The node expansion is done by order of trust followed by cost (the number of new inputs required to compute our target).



\subsection{The source graph}
\label{sec:idea-source}

The \termdef{source graph} is the second knowledge graph. It describes kinds of documents and their official parts. For example a 10-K has Items 1--15 and Item~8 holds the statements, etc. It includes Fifteen kinds and 150 parts based on the official SEC guidelines. Each entry has a description, a verified quote, the time periods a part holds, and its vector embedding.

\subsection{The bridge}
\label{sec:idea-bridge}

The \termdef{bridge} contains the likelihood of finding concepts in any part of a document. It says in which part, statement or note of which kind each quantity is printed and for which time periods. The likelihood is computed using 32,288 random filings from the SEC's Financial Statement and Notes data sets by counting the appearances of each concept in those filings. 

\subsection{Agent Reasoning Loop}
\label{sec:impl-loop}

Figure~\ref{fig:loop} shows the agent loop that answers a question. A language model is called at eight points in the loop, and at each point it makes one narrow choice; everything in between is done by deterministic code. The eight steps are as follows.
\begin{enumerate}[label=(\arabic*)]
  \item \textbf{Understand.} The model first reads the question on its own, before any document is opened, and extracts what is needed to answer it. It returns the targets, that is, the quantities the question asks for, named as the question names them and with their coordinates (company, period and segment). It also returns a formula over these targets, but only when the question spells one out, together with the information given in the question, the unit the answer should be expressed in, and an analysis flag when the question asks for something to be judged rather than only computed.
  \item \textbf{Map.} The words naming each target are then mapped to exactly one entity of the concept graph, or to none if no entity fits. Candidate entities are gathered in three ways: by their names, by the captions that filers print for the corresponding SEC concept, and by vector similarity search. The model is shown the candidates together with their definitions and makes the choice.
  \item \textbf{Choose documents.} The model next chooses which stored documents to read. It does so through tools over the stored document covers, over the periods that each kind of document holds according to the source graph, and over a company store of 16,416 filers and their co-registrants. A prior year that a document holds only as a comparative is ranked below the document's own year, and the tool states this. When only one document fits the question, the choice is made by code without calling the model.
  \item \textbf{Search or derive.} Code then decides whether to search for the target directly or to derive it from other quantities. The decision is based on the bridge's calibrated probability that a filing of this kind, in this industry, prints the entity. If that probability is high, the system searches for the value; if it has been measured as zero, the system tries formulas first; and if it is unknown, the system searches first and derives the value only if the search fails.
  \item \textbf{Routes.} When a target has to be derived, the system collects the relations of the concept graph that produce the entity and solves each of them symbolically for the target. These routes are ranked first by trust and then by the joint probability that all of their inputs will be found. The joint probability is counted in sequence rather than by treating the inputs as independent, because the presence of one quantity in a filing strongly predicts the presence of another: current liabilities is reported by 98.8\% of the filers that report current assets, but by only 3.9\% of the others.
  \item \textbf{Find.} The model then looks for each needed value in the chosen document. It searches first in the parts that the bridge names for that quantity, and only afterwards in the whole document. It reads whole tables together with the text chunks around them, and replies with one of six outcomes: \emph{found}, \emph{found text}, \emph{collect}, \emph{widen}, \emph{search again} or \emph{absent}.
  \item \textbf{Plan.} Code keeps the state of the search, and this state serves as the system's memory from one step to the next. When an input is found, its value is propagated up through every route that it completes. When a route fails because one of its inputs cannot be found, the cheapest of the paused routes is activated in its place.
  \item \textbf{Respond and evaluate.} Finally, every derived value is recomputed in code from its cited inputs, using the equation reloaded from the graph, so that no arithmetic is taken from the model. If the question asks only for a number, the answer is returned without any model call. Otherwise a single model call writes the answer, and every claim in it must quote the evidence it rests on. Code then corrects the shape, unit and sign of the answer.
\end{enumerate}

\begin{figure}[p]
\centering
\fitheight{\fitwidth{\begin{tikzpicture}[
  >=Stealth, font=\scriptsize, node distance=3.5mm,
  step/.style={draw,thick,rounded corners=3pt,align=center,inner sep=4pt,text width=52mm,minimum height=11mm},
  model/.style={step,draw=cModel,fill=cModel!8},
  code/.style={step,draw=cCode,fill=cLight},
  kbnote/.style={draw=cConnector,thick,rounded corners=2pt,fill=cConnector!10,align=left,inner sep=3pt,text width=48mm,font=\tiny},
  arr/.style={->,thick,gray!80},
  back/.style={->,thick,cModel,dashed}
]
\node[model] (und) at (0,0) {\textbf{1 Understand} \textcolor{cModel}{[model]}\\ targets + coordinates, formula over targets, givens, unit, analysis statement};
\node[model,below=of und] (map) {\textbf{2 Map words $\to$ entity} \textcolor{cModel}{[model]}\\ candidates by words, filer captions and meaning; a number or \emph{none}; never cached};
\node[model,below=of map] (docs) {\textbf{3 Choose documents} \textcolor{cModel}{[model, or code when one fits]}\\ tools: search documents, company filers, list};
\node[code,below=of docs] (dec) {\textbf{4 Search or derive?} [code]\\ calibrated chance from the bridge, nearest industry first:\\ high $\to$ search; measured 0 $\to$ formulas; unknown $\to$ search then derive};
\node[model,below=of dec] (routes) {\textbf{5 Routes} \textcolor{cModel}{[model when several]}\\ formulas solved for the target, ranked by trust then joint chance of inputs; ancestors excluded; a definition stated by the model when the graph has none};
\node[model,below=of routes] (find) {\textbf{6 Find} \textcolor{cModel}{[model, rounds $\le 3$]}\\ scope: the parts the bridge names; passages from the retriever; reply: found / found text / collect / widen / search / absent};
\node[code,below=of find] (ground) {\textbf{Ground} [code]\\ quote printed in the passage and row; value and sign from the cell; unit from quoted evidence; kind of number vs kind of quantity};
\node[code,below=of ground] (plan) {\textbf{7 Plan state} [code]\\ settle, propagate up complete routes; switch to a cheaper paused route; abandon inputs no longer needed; evidence kept in full};
\node[model,below=of plan] (resp) {\textbf{8 Respond} \textcolor{cModel}{[model; none for a bare number]}\\ every derived value recomputed from cited inputs via the graph's equation; claims must quote evidence; may ask for more targets};
\node[code,below=of resp] (eval) {\textbf{Evaluate} [code; model only to restate]\\ shape, unit, sign fixed; percentage points for rate differences};
\draw[arr] (und) -- (map); \draw[arr] (map) -- (docs); \draw[arr] (docs) -- (dec); \draw[arr] (dec) -- (routes);
\draw[arr] (dec.east) -- ++(6mm,0) |- node[pos=.25,right,font=\tiny,align=left]{search\\first} (find.east);
\draw[arr] (routes) -- (find); \draw[arr] (find) -- (ground); \draw[arr] (ground) -- (plan); \draw[arr] (plan) -- (resp); \draw[arr] (resp) -- (eval);
\draw[back] (ground.west) -- ++(-5mm,0) |- node[pos=.25,left,font=\tiny,align=right]{rejected:\\ re-ask once} (find.west);
\draw[back] (plan.west) -- ++(-10mm,0) |- node[pos=.3,left,font=\tiny,align=right]{input absent:\\ expand through\\ a formula, or\\ switch route} (routes.west);
\draw[back] (resp.west) -- ++(-15mm,0) |- node[pos=.3,left,font=\tiny,align=right]{needs:\\ new targets\\ ($\le 2$ rounds)} (map.west);

% knowledge-base notes on the right
\node[kbnote,right=12mm of map] (n1) {\textcolor{cEntity}{\textbf{concept graph}}: names, aliases, SEC concept links, filer captions, vectors};
\node[kbnote,right=12mm of docs] (n2) {\textcolor{cDoc}{\textbf{source graph}} + company store: which stored document covers the company and period, as its own year or a comparative};
\node[kbnote,right=12mm of dec] (n3) {\textcolor{cConnector}{\textbf{bridge}}: share of filings of this kind and industry printing the entity; in which part and statement};
\node[kbnote,right=12mm of routes] (n4) {\textcolor{cEntity}{\textbf{concept graph}}: equations and subtotals that produce the entity, solved symbolically; trust: identity $<$ law $<$ empirical};
\node[kbnote,right=12mm of find] (n5) {\textcolor{cConnector}{\textbf{bridge}} $\to$ parts to search first; \textcolor{cDoc}{\textbf{source graph}} $\to$ headings linked to parts; retriever v3 ranks chunks};
\node[kbnote,right=12mm of resp] (n6) {\textcolor{cEntity}{\textbf{concept graph}}: the equation reloaded by connector id for recomputation in code};
\foreach \a/\b in {n1/map,n2/docs,n3/dec,n4/routes,n5/find,n6/resp} {\draw[arr,gray!50] (\a.west) -- (\b.east);}
\end{tikzpicture}}}
\caption{The question-answering loop (version 4). Plum boxes are the points where a model is called, each with one narrow choice; grey boxes are code. Dashed arrows are the loop's backtracking paths. The notes on the right say which knowledge base each step consults.}
\label{fig:loop}
\end{figure}


The rule that runs through the whole loop is \textbf{grounding}. A figure enters an answer only if the code finds the model's quote printed in the cited passage and in the cited row; the value and its sign are then taken from the table cell, not from the model. A unit is accepted when the evidence quoted for it is found in the table heading or the surrounding context, and is otherwise recorded as an assumption. All arithmetic runs in code, with every value normalised to base units and checked for dimensional consistency. What the code deliberately does \emph{not} do is assume the shape of a document: identity tie-outs and range checks are not enforced. When a value disagrees with the graph, the disagreement is put to the model once, and if it remains it is recorded as a flag.

\subsection{Learned where learning is safe}
\label{sec:impl-retrieval}

A hand-ordered search was replaced by \emph{one score per chunk}: thirteen signals (words, synonym, table row label, heading path, meaning, the bridge's share for the part in force, the coordinate's words near and far, table and number counts, length) fitted as a conditional logit~\citep{mcfadden1974} with a ``not printed in this document'' alternative, so it returns a probability over chunks and a probability that the target is absent. The training data is the methodological point: 420 EDGAR filings from two SEC quarters disjoint from the bridge's five, labelled by their own inline-XBRL tags (178,185 values and 23,195 text blocks with the chunks that print them), the benchmark companies excluded, held out by company. No model call and no benchmark question supervises anything. The served probability mass and the bridge's counts were then calibrated on held-out filings~\citep{zadrozny2002}, because both under-predicted in the middle of their range. The principle is general: a piece whose behaviour can be fitted on data the system is never evaluated on is fitted; a judgement about meaning is made by a model every time; a check against the document stays hard and a check against the graph only asks.
